\paragraph{端点三态压缩原理：\texorpdfstring{$\mathcal S_{\star}=\{0,2,-2\}$}{S*={0,2,-2}} 与模 \texorpdfstring{$4$}{4} 三通道}
为把命题 \ref{prop:half-angle-z4-residue} 的 \(\ZZ/4\) 残差进一步压缩为“有限状态”模板，引入端点三态集合。

\begin{definition}[端点三态集合与端点态函数]\label{def:endpoint-tristate-set}
在端点 \(w=-1\)（等价于 \(u=-1\)，且 \(S=0\)）处，记
\[
\boxed{\ \mathcal S_{\star}:=\{0,2,-2\}\ }.
\]
并称映射
\[
d\ \longmapsto\ C_d(0)\in\mathcal S_{\star}\qquad(d\ge 1)
\]
为幂伸缩指数 \(d\) 的\textbf{端点态函数}。
\end{definition}

\leanverified{paper\_endpoint\_tristate\_compression}
\begin{corollary}[端点三态压缩原理（\texorpdfstring{$\bmod 4$}{mod 4} 三通道退化）]\label{cor:endpoint-tristate-compression}
对一切整数 \(d\ge 1\)，端点取值满足
\begin{equation}\label{eq:endpoint-tristate-values}
\boxed{\;
C_d(0)=\mathrm{i}^{d}+\mathrm{i}^{-d}=2\cos(d\pi/2)\in\mathcal S_{\star},\qquad
C_d(0)=
\begin{cases}
0, & d\ \text{为奇数},\\
2, & d\equiv 0\ (\mathrm{mod}\ 4),\\
-2,& d\equiv 2\ (\mathrm{mod}\ 4).
\end{cases}\;}
\end{equation}
因此，任意由幂替换半群 \(\{u\mapsto u^d\}_{d\ge 1}\) 诱导的完成化动力学 \(S\mapsto C_d(S)\) 在端点 \(S=0\) 的 \(0\)-喷射上必然退化为一个模 \(4\) 的三态系统。
\end{corollary}
\begin{proof}
式 \eqref{eq:endpoint-tristate-values} 即命题 \ref{prop:half-angle-z4-residue} 的端点恒等式 \eqref{eq:chebyshev-endpoint-values} 在三类剩余 \(d\bmod 4\) 下的展开写法。证毕。
\end{proof}

\paragraph{Witt--Chebyshev 端点喷射分解：primitive 的三态模板}
令 \(p_n(u)\in\ZZ[u]\) 为 primitive 多项式，并在完成化坐标 \(u=r^2,\ S=r+r^{-1}\) 下定义
\(\widehat p_n(S):=r^{-n}p_n(r^2)\in\ZZ[S]\)；
同时对任意 \(d\ge 1\) 记 \(C_d(S):=r^d+r^{-d}\in\ZZ[S]\)。
则有严格倍角律 \(r^{-dn}p_n(r^{2d})=\widehat p_n(C_d(S))\)（见推论 \ref{cor:primitive-completion-hatp}）。
据此引入端点三态值
\begin{equation}\label{eq:primitive-endpoint-tristate-values}
\boxed{\;
A_n:=\widehat p_n(0),\qquad
B_n^{+}:=\widehat p_n(2),\qquad
B_n^{-}:=\widehat p_n(-2)\; }.
\end{equation}

\begin{proposition}[primitive 的端点三态退化：\texorpdfstring{$\widehat p_n(C_d(0))$}{hat p_n(C_d(0))} 的模 \texorpdfstring{$4$}{4} 通道分裂]\label{prop:primitive-endpoint-tristate-degeneration}
\leanverified{paper\_primitive\_endpoint\_tristate\_degeneration}
对任意 \(n\ge 1\) 与 \(d\ge 1\)，在端点 \(S=0\) 有严格退化
\begin{equation}\label{eq:primitive-endpoint-tristate-degeneration}
\boxed{\;
\widehat p_n\!\bigl(C_d(0)\bigr)=
\begin{cases}
A_n, & d \text{ 奇},\\
B_n^{+}, & d\equiv 0\ (\mathrm{mod}\ 4),\\
B_n^{-}, & d\equiv 2\ (\mathrm{mod}\ 4).
\end{cases}\;}
\end{equation}
并且端点两侧 \(\pm 2\) 与 \(u=1\) 的关系给出
\begin{equation}\label{eq:primitive-endpoint-Bpm-identification}
\boxed{\ B_n^{+}=p_n(1),\qquad B_n^{-}=(-1)^n p_n(1)\ }.
\end{equation}
\end{proposition}
\begin{proof}
由推论 \ref{cor:endpoint-tristate-compression}，端点态 \(C_d(0)\) 仅取三值 \(\{0,2,-2\}\)，且其取值由 \(d\bmod 4\) 决定；代入 \(\widehat p_n\) 即得 \eqref{eq:primitive-endpoint-tristate-degeneration}。
\par
另一方面，由定义 \(\widehat p_n(S)=r^{-n}p_n(r^2)\)，当 \(S=2\) 时取 \(r=1\) 得 \(\widehat p_n(2)=p_n(1)\)。
当 \(S=-2\) 时取 \(r=-1\) 得 \(\widehat p_n(-2)=(-1)^{-n}p_n(1)=(-1)^n p_n(1)\)。
从而得到 \eqref{eq:primitive-endpoint-Bpm-identification}。证毕。
\end{proof}

\paragraph{端点 WITT 分解模板：三通道权重核与对数主尺度}
对在包含区间 \([-2,2]\) 的某开邻域上解析的函数 \(g\)，定义截断的完成化幂链算子
\begin{equation}\label{eq:completed-witt-pow-operator-trunc}
\boxed{\;
(W_{\mathrm{pow}}^{\le M}g)(S):=\sum_{d=1}^{M}\frac{1}{d}\,g\!\bigl(C_d(S)\bigr)\qquad(M\ge 1)\; }.
\end{equation}
则端点 \(S=0\) 的 \(0\)-喷射满足严格的模类通道分解：

\begin{proposition}[端点三通道分解：截断幂链在 \texorpdfstring{$S=0$}{S=0} 的模 \texorpdfstring{$4$}{4} 分组恒等式]\label{prop:endpoint-mod4-channel-decomposition}
令
\[
H_M^{(1)}:=\sum_{\substack{1\le d\le M\\ d\ \mathrm{odd}}}\frac1d,\qquad
H_M^{(0)}:=\sum_{\substack{1\le d\le M\\ d\equiv 0\ (\mathrm{mod}\ 4)}}\frac1d,\qquad
H_M^{(2)}:=\sum_{\substack{1\le d\le M\\ d\equiv 2\ (\mathrm{mod}\ 4)}}\frac1d.
\]
则对任意解析 \(g\) 有恒等式
\begin{equation}\label{eq:endpoint-mod4-channel-decomposition}
\boxed{\;
(W_{\mathrm{pow}}^{\le M}g)(0)
=
H_M^{(1)}\,g(0)
+H_M^{(0)}\,g(2)
+H_M^{(2)}\,g(-2)\; }.
\end{equation}
因此端点处的幂链叠加在 \(0\)-喷射层面退化为三个互不耦合的模类通道，权重完全由 \(\sum_{d\equiv a\ (\mathrm{mod}\ 4)}\frac1d\) 型核决定。
\end{proposition}
\begin{proof}
由定义 \eqref{eq:completed-witt-pow-operator-trunc} 与端点三态压缩 \eqref{eq:endpoint-tristate-values}，
对每一项 \(g(C_d(0))\) 按 \(d\bmod 4\) 分类并合并即可得到 \eqref{eq:endpoint-mod4-channel-decomposition}。证毕。
\end{proof}
\leanverified{paper\_endpoint\_mod4\_channel\_decomposition}

\begin{corollary}[奇通道的结构性空化（端点 \(0\)-喷射）]\label{cor:endpoint-odd-channel-vanish}
若解析观测 \(g\) 满足 \(g(0)=0\)，则对任意 \(M\ge 1\) 有
\[
\boxed{\;
(W_{\mathrm{pow}}^{\le M}g)(0)=H_M^{(0)}\,g(2)+H_M^{(2)}\,g(-2)\; }.
\]
因此端点三通道中由奇 \(d\) 通道所携带的 \(g(0)\) 分量完全消失；端点边界层的 \(0\)-喷射信息仅由两条偶通道 \(S=\pm 2\) 的对比给出。
\end{corollary}
\begin{proof}
由 \eqref{eq:endpoint-mod4-channel-decomposition} 代入 \(g(0)=0\) 立得。证毕。
\end{proof}
\leanverified{paper\_endpoint\_odd\_channel\_vanish}

\begin{proposition}[端点三通道的 Dirichlet 正则化权核与闭式表示]\label{prop:endpoint-dirichlet-symbol-weights}
对 \(\sigma>0\) 定义端点正则化幂链算子
\begin{equation}\label{eq:endpoint-regularized-channel-operator}
\boxed{\;
(W_{\mathrm{pow}}^{(\sigma)}g)(0)
:=\sum_{d\ge 1}\frac{1}{d^{1+\sigma}}\,g\!\bigl(C_d(0)\bigr)\; }.
\end{equation}
则其端点三通道分解为
\begin{equation}\label{eq:endpoint-regularized-channel-decomposition}
\boxed{\;
(W_{\mathrm{pow}}^{(\sigma)}g)(0)
=\mathsf H_{\mathrm{odd}}(\sigma)\,g(0)
+\mathsf H_{0}(\sigma)\,g(2)
+\mathsf H_{2}(\sigma)\,g(-2)\; }.
\end{equation}
其中三通道权核具有闭式 Dirichlet 符号表示
\begin{align}
\mathsf H_{\mathrm{odd}}(\sigma)
&:=\sum_{\substack{d\ge 1\\ d\ \mathrm{odd}}}\frac{1}{d^{1+\sigma}}
=\bigl(1-2^{-(1+\sigma)}\bigr)\zeta(1+\sigma),\label{eq:endpoint-Hodd-sigma}\\
\mathsf H_{0}(\sigma)
&:=\sum_{d\equiv 0\ (\mathrm{mod}\ 4)}\frac{1}{d^{1+\sigma}}
=4^{-(1+\sigma)}\zeta(1+\sigma),\label{eq:endpoint-H0-sigma}\\
\mathsf H_{2}(\sigma)
&:=\sum_{d\equiv 2\ (\mathrm{mod}\ 4)}\frac{1}{d^{1+\sigma}}
=\bigl(2^{-(1+\sigma)}-4^{-(1+\sigma)}\bigr)\zeta(1+\sigma).\label{eq:endpoint-H2-sigma}
\end{align}
\end{proposition}
\begin{proof}
由端点三态压缩 \eqref{eq:endpoint-tristate-values}，\(C_d(0)\) 仅依赖 \(d\bmod 4\) 并取值于 \(\{0,2,-2\}\)，故
\[
(W_{\mathrm{pow}}^{(\sigma)}g)(0)
=\sum_{\substack{d\ge 1\\ d\ \mathrm{odd}}}\frac{1}{d^{1+\sigma}}\,g(0)
+\sum_{d\equiv 0(4)}\frac{1}{d^{1+\sigma}}\,g(2)
+\sum_{d\equiv 2(4)}\frac{1}{d^{1+\sigma}}\,g(-2),
\]
即得 \eqref{eq:endpoint-regularized-channel-decomposition}。
对闭式表示，注意到
\[
\sum_{\substack{d\ge 1\\ d\ \mathrm{odd}}}\frac{1}{d^{1+\sigma}}
\;=\;\sum_{d\ge 1}\frac{1}{d^{1+\sigma}}-\sum_{k\ge 1}\frac{1}{(2k)^{1+\sigma}}
\;=\;\bigl(1-2^{-(1+\sigma)}\bigr)\zeta(1+\sigma),
\]
并且
\[
\sum_{d\equiv 0(4)}\frac{1}{d^{1+\sigma}}
=\sum_{k\ge 1}\frac{1}{(4k)^{1+\sigma}}
=4^{-(1+\sigma)}\zeta(1+\sigma),
\]
以及
\(
\sum_{d\equiv 2(4)}d^{-(1+\sigma)}
=\sum_{k\ge 1}(2k)^{-(1+\sigma)}-\sum_{k\ge 1}(4k)^{-(1+\sigma)}
\)
给出 \eqref{eq:endpoint-H2-sigma}。证毕。
\end{proof}
\leanverified{paper\_endpoint\_dirichlet\_symbol\_weights}

\paragraph{端点相位的四通道提升：\texorpdfstring{$\ZZ/4$}{Z/4} 分支记忆在完成化层的最小闭包}
在完成化坐标 \(u=r^{2},\ S=r+r^{-1}\) 下，映射 \(r\mapsto r^{-1}\) 的反演不变子环严格下降为 \(\ZZ[S]\)（命题 \ref{prop:completion-subring-unique-angle}），从而端点 \(u=-1\) 的可见信息在 \(0\)-喷射层面退化为三态 \(\mathcal S_\star=\{0,\pm 2\}\)（推论 \ref{cor:endpoint-tristate-compression}）。
为在完成化层捕获端点的 \(\ZZ/4\) 分支记忆，引入最小的反演反变伴随坐标
\begin{equation}\label{eq:endpoint-delta-anti-invariant}
\boxed{\ \delta:=r-r^{-1}\ }.
\end{equation}
并对任意 \(d\ge 1\) 定义“奇”Chebyshev--Adams 伴随多项式族
\begin{equation}\label{eq:endpoint-odd-chebyshev-adams}
\boxed{\ D_d(S,\delta):=r^{d}-r^{-d}\ }.
\end{equation}

\begin{remark}[反演反变扇区的主生成元与最小闭包]\label{rem:endpoint-inv-antiinv-min-closure}
令
\[
\mathcal A^{-}:=\{\,f(r)\in \ZZ[r,r^{-1}]:\ f(r^{-1})=-f(r)\,\}
\]
为反演作用 \(r\mapsto r^{-1}\) 下的反变部分。
则有
\[
\boxed{\ \mathcal A^{-}=\delta\,\ZZ[S]\ },\qquad S=r+r^{-1},\ \delta=r-r^{-1}.
\]
因此在完成化层面，将 \(\delta\) 作为额外纤维坐标即可生成全部反变观测，并在端点分支 \(r=\pm\mathrm{i}\) 上以 \(\delta=\pm 2\mathrm{i}\) 读出分支记忆；这给出从三态 \(\mathcal S_\star\) 到完整 \(\ZZ/4\) 通道的最小闭包。
并且 \(\delta^2=S^2-4\)，故将可观测环从 \(\ZZ[S]\) 提升到满足关系 \(\delta^2=S^2-4\) 的最小二次闭包，可在不引入额外高维自由度的前提下恢复端点分支信息。
\end{remark}

\begin{proposition}[反变扇区的 Frobenius-lift：\texorpdfstring{$(S,\delta)\mapsto(C_p(S),D_p(S,\delta))$}{(S,δ)->(C_p(S),D_p(S,δ))}]\label{prop:endpoint-antiinv-frobenius-lift}
\leanverified{paper\_endpoint\_antiinv\_frobenius\_lift}
令最小二次闭包
\[
\mathcal{A}_{\min}:=\ZZ[S,\delta]\big/\bigl(\delta^2-(S^2-4)\bigr).
\]
对任意素数 \(p\)，定义映射 \(\varphi_p:\mathcal{A}_{\min}\to\mathcal{A}_{\min}\) 为
\[
\varphi_p(S):=C_p(S),
\qquad
\varphi_p(\delta):=D_p(S,\delta).
\]
则 \(\varphi_p\) 为良定义的环自同态，并在模 \(p\) 层面与 Frobenius 对齐：
\[
\varphi_p(S)\equiv S^p\pmod p,
\qquad
\varphi_p(\delta)\equiv \delta^p\pmod p.
\]
因此 \(\varphi_p\) 在不变子环 \(\ZZ[S]\) 上与命题 \ref{prop:chebyshev-frobenius-pderivation-integrality} 的 Chebyshev--Frobenius 提升一致，而在反变纤维上给出与端点 \(\ZZ/4\) 分支记忆相容的最小扩张 Frobenius-lift。
\end{proposition}
\begin{proof}
在 \(\ZZ[r,r^{-1}]\) 中有恒等式 \(C_p(S)=r^p+r^{-p}\)、\(D_p(S,\delta)=r^p-r^{-p}\)，从而
\[
D_p(S,\delta)^2
=\bigl(r^p-r^{-p}\bigr)^2
=\bigl(r^p+r^{-p}\bigr)^2-4
=C_p(S)^2-4,
\]
故关系 \(\delta^2=S^2-4\) 在 \(\varphi_p\) 下保持，映射良定义。
模 \(p\) 的同余由二项式系数的 \(p\)-整除性给出：
\[
(r+r^{-1})^p\equiv r^p+r^{-p}\ (\mathrm{mod}\ p),
\qquad
(r-r^{-1})^p\equiv r^p-r^{-p}\ (\mathrm{mod}\ p),
\]
从而 \(C_p(S)\equiv S^p\) 且 \(D_p(S,\delta)\equiv \delta^p\)。
对 \(f\in\ZZ[S]\) 的限制即 \(f(S)\mapsto f(C_p(S))\)，与命题 \ref{prop:chebyshev-frobenius-pderivation-integrality} 的 Frobenius-lift 恒等式一致。证毕。
\end{proof}

\begin{proposition}[奇 Chebyshev--Adams 伴随恒等式]\label{prop:endpoint-odd-chebyshev-identity}
\leanverified{paper\_endpoint\_odd\_chebyshev\_identity}
对任意 \(d\ge 1\)，有恒等式
\begin{equation}\label{eq:endpoint-odd-chebyshev-identity}
\boxed{\;
D_d(S,\delta)=\delta\,U_{d-1}(S/2)\; },
\end{equation}
其中 \(U_{d-1}\) 为第二类 Chebyshev 多项式。
特别地，在端点 \(u=-1\) 的半角分支 \(r=\pm\mathrm{i}\) 上有 \(S=0\)、\(\delta=\pm 2\mathrm{i}\)，且
\begin{equation}\label{eq:endpoint-odd-chebyshev-ratio}
\boxed{\;
\frac{D_d(0,\delta)}{\delta}
=\frac{r^{d}-r^{-d}}{r-r^{-1}}
=U_{d-1}(0)
=\sin(d\pi/2)\; }.
\end{equation}
\end{proposition}
\begin{proof}
令 \(E_d:=r^{d}-r^{-d}\)。则由 \(r+r^{-1}=S\) 直接验证递推
\(
E_{d+1}=SE_d-E_{d-1}
\)
并有初值 \(E_0=0,E_1=\delta\)。
因此 \(E_d/\delta\) 满足与 \(U_{d-1}(S/2)\) 相同的递推与初值（\(U_0=1,U_1=S\)），从而得到 \eqref{eq:endpoint-odd-chebyshev-identity}。
\par
取端点 \(S=0\) 并用恒等式 \(U_{d-1}(\cos\phi)=\sin(d\phi)/\sin\phi\)（令 \(\phi=\pi/2\)）得 \(U_{d-1}(0)=\sin(d\pi/2)\)，从而得到 \eqref{eq:endpoint-odd-chebyshev-ratio}。证毕。
\end{proof}

\begin{theorem}[端点四通道闭包与 \texorpdfstring{$\chi_{4}$}{chi4} 读出]\label{thm:endpoint-four-channel-chi4}
\leanverified{paper\_endpoint\_four\_channel\_chi4}
定义端点通道指标
\begin{equation}\label{eq:endpoint-epsilon-def}
\boxed{\ \varepsilon(d):=\sin(d\pi/2)\in\{-1,0,1\}\qquad(d\ge 1)\ }.
\end{equation}
则 \(\varepsilon(d)=0\) 当且仅当 \(d\) 为偶数；当 \(d\) 为奇数时，
\[
\varepsilon(d)=
\begin{cases}
1, & d\equiv 1\ (\mathrm{mod}\ 4),\\
-1,& d\equiv 3\ (\mathrm{mod}\ 4).
\end{cases}
\]
因此 \(\varepsilon\) 给出模 \(4\) 的完整奇类分裂。
令 \(\chi_{4}\) 为模 \(4\) 的非平凡 Dirichlet 角色
\[
\boxed{\;
\chi_{4}(d)=
\begin{cases}
0, & d\ \text{为偶数},\\
1, & d\equiv 1\ (\mathrm{mod}\ 4),\\
-1,& d\equiv 3\ (\mathrm{mod}\ 4),
\end{cases}\;}
\]
则对一切 \(d\ge 1\) 有 \(\chi_{4}(d)=\varepsilon(d)\)。
该结论表明：仅靠完成化偶部分 \(C_d\) 只能得到端点三态，而加入最小的反演反变纤维坐标 \(\delta\) 后，端点相位的分支记忆自然闭包为完整的 \(\ZZ/4\) 通道。
\end{theorem}
\begin{proof}
由 \(\sin(d\pi/2)\) 的取值仅依赖 \(d\bmod 4\)，并直接列举四类剩余即可得到上述分裂；与 \(\chi_4\) 的逐类定义一致，故 \(\chi_4(d)=\varepsilon(d)\)。证毕。
\end{proof}

\paragraph{端点有限维不变子表示：\texorpdfstring{$\zeta$}{zeta} 与 \texorpdfstring{$L(\chi_{4})$}{L(chi4)} 的二重符号}
对 \(\sigma>0\)，在完成化坐标上引入正则化的 Chebyshev--Witt 幂链算子
\begin{equation}\label{eq:endpoint-Wsigma-even}
\boxed{\;
(\mathcal W_{\sigma}F)(S):=\sum_{m\ge 1}\frac{1}{m^{1+\sigma}}\,F\!\bigl(C_m(S)\bigr)\; },
\qquad S\in[-2,2],\ F\in\CC[S],
\end{equation}
及其在反变纤维扩展 \((S,\delta)\) 上的伴随作用
\begin{equation}\label{eq:endpoint-Wsigma-odd}
\boxed{\;
(\mathcal W_{\sigma}^{\delta}G)(S,\delta):=\sum_{m\ge 1}\frac{1}{m^{1+\sigma}}\,G\!\bigl(C_m(S),D_m(S,\delta)\bigr)\; },
\qquad S\in[-2,2],\ G\in\CC[S,\delta].
\end{equation}
并定义三个偶端点线性泛函与一个反变端点线性泛函：
\[
\mathcal E_{0}(F):=F(0),\qquad \mathcal E_{+}(F):=F(2),\qquad \mathcal E_{-}(F):=F(-2),
\]
\begin{equation}\label{eq:endpoint-odd-functional-O}
\boxed{\ \mathcal O(G):=\frac{1}{2\mathrm{i}}\Bigl(G(0,2\mathrm{i})-G(0,-2\mathrm{i})\Bigr)\ }.
\end{equation}

\begin{theorem}[端点符号分解：\texorpdfstring{$\zeta$}{zeta}--\texorpdfstring{$L(\chi_{4})$}{L(chi4)} 的正交二元组]\label{thm:endpoint-zeta-Lchi4-symbol-pair}
\leanverified{paper\_endpoint\_zeta\_lchi4\_symbol\_pair}
对任意 \(\sigma>0\) 与 \(F\in\CC[S]\)，有端点偶观测分解
\begin{equation}\label{eq:endpoint-even-symbol-decomposition-alpha}
\boxed{\;
\mathcal E_{0}(\mathcal W_{\sigma}F)
=\alpha_{0}(\sigma)\,\mathcal E_{0}(F)+\alpha_{+}(\sigma)\,\mathcal E_{+}(F)+\alpha_{-}(\sigma)\,\mathcal E_{-}(F)\; },
\end{equation}
其中
\[
\alpha_{0}(\sigma):=\sum_{\substack{m\ge 1\\ m\ \mathrm{odd}}}\frac{1}{m^{1+\sigma}},\qquad
\alpha_{+}(\sigma):=\sum_{m\equiv 0(4)}\frac{1}{m^{1+\sigma}},\qquad
\alpha_{-}(\sigma):=\sum_{m\equiv 2(4)}\frac{1}{m^{1+\sigma}}.
\]
并且 \(\alpha_{0},\alpha_{+},\alpha_{-}\) 具有闭式 \(\zeta\)-符号表示（命题 \ref{prop:endpoint-dirichlet-symbol-weights} 的 \eqref{eq:endpoint-Hodd-sigma}--\eqref{eq:endpoint-H2-sigma}）。
\par
同时，对任意 \(G\in\CC[S,\delta]\) 有端点反变观测的本征分解
\begin{equation}\label{eq:endpoint-odd-symbol-eigen}
\boxed{\;
\mathcal O(\mathcal W_{\sigma}^{\delta}G)
=L(1+\sigma,\chi_{4})\cdot \mathcal O(G)\; },
\end{equation}
其中
\(
L(1+\sigma,\chi_{4}):=\sum_{m\ge 1}\chi_{4}(m)m^{-(1+\sigma)}
\)
为模 \(4\) 角色的 Dirichlet \(L\)-级数。
换言之，端点反变观测 \(\mathcal O\) 在 \(\mathcal W_{\sigma}^{\delta}\) 下构成一维不变子表示，其符号恰为 \(L(1+\sigma,\chi_{4})\)；并与偶通道的 \(\zeta\)-型权核正交解耦。
\end{theorem}
\begin{proof}
偶观测 \eqref{eq:endpoint-even-symbol-decomposition-alpha} 由端点三态压缩 \eqref{eq:endpoint-tristate-values} 将求和按 \(m\bmod 4\) 分类得到；其闭式表达见命题 \ref{prop:endpoint-dirichlet-symbol-weights}。
\par
对反变观测，利用 \eqref{eq:endpoint-odd-functional-O} 与 \eqref{eq:endpoint-Wsigma-odd} 得
\[
\mathcal O(\mathcal W_{\sigma}^{\delta}G)
=\frac{1}{2\mathrm{i}}\sum_{m\ge 1}\frac{1}{m^{1+\sigma}}
\Bigl(G(C_m(0),D_m(0,2\mathrm{i}))-G(C_m(0),D_m(0,-2\mathrm{i}))\Bigr).
\]
若 \(m\) 为偶数，则由 \eqref{eq:endpoint-epsilon-def} 有 \(\varepsilon(m)=0\)，从而由命题 \ref{prop:endpoint-odd-chebyshev-identity} 得 \(D_m(0,\pm 2\mathrm{i})=0\)，对应项在差分中严格抵消。
若 \(m\) 为奇数，则 \(C_m(0)=0\) 且 \(\varepsilon(m)=\pm 1\)，并由 \eqref{eq:endpoint-odd-chebyshev-ratio} 得
\(
D_m(0,\pm 2\mathrm{i})=\pm 2\mathrm{i}\,\varepsilon(m)
\)。
于是
\[
G(0,2\mathrm{i}\varepsilon(m))-G(0,-2\mathrm{i}\varepsilon(m))
=2\mathrm{i}\,\varepsilon(m)\,\mathcal O(G).
\]
代回即得
\[
\mathcal O(\mathcal W_{\sigma}^{\delta}G)
=\left(\sum_{m\ge 1}\frac{\varepsilon(m)}{m^{1+\sigma}}\right)\mathcal O(G)
=\left(\sum_{m\ge 1}\frac{\chi_{4}(m)}{m^{1+\sigma}}\right)\mathcal O(G)
=L(1+\sigma,\chi_{4})\,\mathcal O(G),
\]
其中第二个等号由定理 \ref{thm:endpoint-four-channel-chi4}（\(\varepsilon=\chi_4\)）给出。证毕。
\end{proof}

\begin{proposition}[端点偶三态的不变子表示与显式矩阵化]\label{prop:endpoint-even-tristate-matrix}
\leanverified{paper\_endpoint\_even\_tristate\_matrix}
沿用 \eqref{eq:endpoint-Wsigma-even}--\eqref{eq:endpoint-odd-symbol-eigen} 的记号，并令
\[
v(F):=\bigl(\mathcal E_{0}(F),\mathcal E_{+}(F),\mathcal E_{-}(F)\bigr)^{\mathsf T}\in\CC^{3}.
\]
则对任意 \(\sigma>0\) 与任意 \(F\in\CC[S]\)，端点偶三态在 \(\mathcal W_\sigma\) 下封闭，并满足
\begin{equation}\label{eq:endpoint-even-tristate-matrix-action}
\boxed{\;
v(\mathcal W_{\sigma}F)=M_{\mathrm{even}}(\sigma)\,v(F)\;}
\end{equation}
其中表示矩阵可写成闭式
\begin{equation}\label{eq:endpoint-even-tristate-matrix-closed}
\boxed{\;
M_{\mathrm{even}}(\sigma)
=\zeta(1+\sigma)
\begin{pmatrix}
1-2^{-(1+\sigma)} & 4^{-(1+\sigma)} & 2^{-(1+\sigma)}-4^{-(1+\sigma)}\\
0&1&0\\
0&2^{-(1+\sigma)}&1-2^{-(1+\sigma)}
\end{pmatrix}. \;}
\end{equation}
\end{proposition}

\begin{proof}
第一行由定理 \ref{thm:endpoint-zeta-Lchi4-symbol-pair} 的端点偶观测分解 \eqref{eq:endpoint-even-symbol-decomposition-alpha} 与 \(\alpha\)-权核的闭式表达给出。
\par
对 \(\mathcal E_{+}\) 与 \(\mathcal E_{-}\)，由完成化迹坐标的倍角恒等式 \(C_m(2)=2\) 与 \(C_m(-2)=2(-1)^m\) 得
\[
\mathcal E_{+}(\mathcal W_\sigma F)
=\sum_{m\ge 1}\frac{1}{m^{1+\sigma}}\,F\!\bigl(C_m(2)\bigr)
=\zeta(1+\sigma)\,\mathcal E_{+}(F),
\]
\[
\mathcal E_{-}(\mathcal W_\sigma F)
=\sum_{\substack{m\ge 1\\ m\ \mathrm{even}}}\frac{1}{m^{1+\sigma}}\,F(2)
\;+\;
\sum_{\substack{m\ge 1\\ m\ \mathrm{odd}}}\frac{1}{m^{1+\sigma}}\,F(-2)
=2^{-(1+\sigma)}\zeta(1+\sigma)\,\mathcal E_{+}(F)
\;+\;\bigl(1-2^{-(1+\sigma)}\bigr)\zeta(1+\sigma)\,\mathcal E_{-}(F).
\]
将上述三式写成矩阵形式即得 \eqref{eq:endpoint-even-tristate-matrix-action}--\eqref{eq:endpoint-even-tristate-matrix-closed}。证毕。
\end{proof}

\begin{theorem}[端点偶通道的普适极与有限部（\texorpdfstring{$\sigma\downarrow 0$}{sigma->0}）]\label{thm:endpoint-even-finite-part-renormalization}
\leanverified{paper\_endpoint\_even\_finite\_part\_renormalization}
令
\[
\mu:=\left(\frac12,\frac14,\frac14\right),\qquad
\langle \mu,(v_0,v_+,v_-)\rangle:=\frac12 v_0+\frac14 v_++\frac14 v_-.
\]
对任意 \(F\in\CC[S]\)，定义端点偶通道的重整化有限部泛函
\begin{equation}\label{eq:endpoint-even-renormalized-finite-part-def}
\boxed{\;
\mathcal R_{\mathrm{even}}(F)
:=\lim_{\sigma\downarrow 0}\left(
\mathcal E_{0}(\mathcal W_\sigma F)
-\frac{1}{\sigma}\,\langle \mu,v(F)\rangle
\right)\;}
\end{equation}
（若极限存在）。
则该极限存在，并给出显式有限部
\begin{equation}\label{eq:endpoint-even-renormalized-finite-part-closed}
\boxed{\;
\mathcal R_{\mathrm{even}}(F)
=\left(\frac{\gamma}{2}+\frac{\log 2}{2}\right)\mathcal E_{0}(F)
\;+\;\left(\frac{\gamma}{4}-\frac{\log 2}{2}\right)\mathcal E_{+}(F)
\;+\;\frac{\gamma}{4}\mathcal E_{-}(F)\; }.
\end{equation}
其中 \(\gamma\) 为 Euler 常数。
\end{theorem}

\begin{proof}
由定理 \ref{thm:endpoint-zeta-Lchi4-symbol-pair} 与 \eqref{eq:endpoint-even-symbol-decomposition-alpha}，
\[
\mathcal E_0(\mathcal W_\sigma F)
=\alpha_0(\sigma)\mathcal E_0(F)+\alpha_+(\sigma)\mathcal E_+(F)+\alpha_-(\sigma)\mathcal E_-(F),
\]
并且（命题 \ref{prop:endpoint-dirichlet-symbol-weights}）
\[
\alpha_0(\sigma)=\bigl(1-2^{-(1+\sigma)}\bigr)\zeta(1+\sigma),\quad
\alpha_+(\sigma)=4^{-(1+\sigma)}\zeta(1+\sigma),\quad
\alpha_-(\sigma)=\bigl(2^{-(1+\sigma)}-4^{-(1+\sigma)}\bigr)\zeta(1+\sigma).
\]
利用 Laurent 展开 \(\zeta(1+\sigma)=\sigma^{-1}+\gamma+O(\sigma)\)（\(\sigma\downarrow 0\)）以及
\[
2^{-(1+\sigma)}=\frac12\bigl(1-\sigma\log 2+O(\sigma^2)\bigr),\qquad
4^{-(1+\sigma)}=\frac14\bigl(1-2\sigma\log 2+O(\sigma^2)\bigr),
\]
可得
\[
\alpha_0(\sigma)=\frac{1}{2\sigma}+\left(\frac{\gamma}{2}+\frac{\log 2}{2}\right)+O(\sigma),\quad
\alpha_+(\sigma)=\frac{1}{4\sigma}+\left(\frac{\gamma}{4}-\frac{\log 2}{2}\right)+O(\sigma),\quad
\alpha_-(\sigma)=\frac{1}{4\sigma}+\frac{\gamma}{4}+O(\sigma).
\]
将其代回并按 \(\sigma^{-1}\) 极点项与常数项分离，即得到 \eqref{eq:endpoint-even-renormalized-finite-part-def} 的极限存在并等于 \eqref{eq:endpoint-even-renormalized-finite-part-closed}。证毕。
\end{proof}

\begin{corollary}[反变端点通道的解析性与 \texorpdfstring{$\chi_4$}{chi4} 常数输出]\label{cor:endpoint-odd-analytic-at-zero}
\leanverified{paper\_endpoint\_odd\_analytic\_at\_zero}
对任意 \(G\in\CC[S,\delta]\)，有
\begin{equation}\label{eq:endpoint-odd-limit-L1-chi4}
\boxed{\;
\lim_{\sigma\downarrow 0}\mathcal O(\mathcal W_{\sigma}^{\delta}G)
=L(1,\chi_{4})\cdot \mathcal O(G)
=\frac{\pi}{4}\,\mathcal O(G)\; }.
\end{equation}
\end{corollary}

\begin{proof}
由定理 \ref{thm:endpoint-zeta-Lchi4-symbol-pair} 的本征分解 \eqref{eq:endpoint-odd-symbol-eigen}，\(\mathcal O(\mathcal W_{\sigma}^{\delta}G)=L(1+\sigma,\chi_4)\mathcal O(G)\)。
由于 \(L(1+\sigma,\chi_4)\) 在 \(\sigma=0\) 处解析，取极限得第一等号；经典恒等式 \(L(1,\chi_4)=\pi/4\) 给出第二等号。证毕。
\end{proof}

\begin{remark}[端点重整化的二元残余：普适极 \(\oplus\) 离散本征通道]\label{rem:endpoint-renormalization-universal-pole-discrete-hair}
定理 \ref{thm:endpoint-even-finite-part-renormalization} 表明：端点偶通道的全部发散仅沿
\(\langle\mu,v(F)\rangle\) 的 rank-one 方向出现，并由 \(\zeta(1+\sigma)\) 在 \(1\) 处的极点完全控制；
其有限部由 Euler 常数与 \(\log 2\) 的显式组合给出。
与之正交地，推论 \ref{cor:endpoint-odd-analytic-at-zero} 表明反变端点观测在 \(\sigma\downarrow 0\) 处无极点，且其残余压缩为纯离散的 \(\chi_4\) 常数通道 \(L(1,\chi_4)=\pi/4\)。
因此，在端点符号口径下，重整化后的非紧性可被压缩为“普适 \(\zeta\)-极 + 离散 \(\chi_4\) 本征残余”的二元正交分解。
\end{remark}

\begin{remark}[端点对数主尺度与相位精度解耦：\texorpdfstring{$u=-1$}{u=-1} 下偶幂通道的计数控制]\label{rem:endpoint-logscale-parity-decouple}
在 twist 层 \(u=-1\) 上，偶幂满足 \(u^{2k}=1\)，故任何 Witt 型伸缩叠加
\(
\sum_{m\ge 1}\frac1m f(u^m)
\)
都包含一个由偶幂子链 \(\sum_{k\ge 1}\frac{1}{2k}f(1)\) 诱导的对数主尺度。
对 primitive 输入 \(f=p_n\) 而言，该主尺度的系数为
\(
f(1)=p_n(1)=B_n^{+}
\)
（见 \eqref{eq:primitive-endpoint-Bpm-identification}），其与端点差分证书
\(
A_n=\widehat p_n(0)
\)
所编码的半角分支信息在定义上分离。
特别地，当 \(n\) 为奇数时 \(\widehat p_n\) 为奇多项式（推论 \ref{cor:primitive-completion-hatp}），从而 \(A_n=\widehat p_n(0)=0\)；而当 \(n\ge 3\) 为奇数时还满足 \(p_n(-1)=0\)（推论 \ref{cor:primitive-odd-vanish-u-minus1}），端点奇偶差分证书在奇长 primitive 上完全消失。
\end{remark}

\paragraph{Frobenius--Adams 整数层相位证书：\texorpdfstring{$\Delta_{p,k}$}{Delta} 在端点三态上的分层审计}
定义 \ref{def:witt-frobenius-defect} 给出 Frobenius 缺陷多项式
\(
\Delta_{p,k}(u):=a_{p^k}(u)-a_{p^{k-1}}(u^p)\in\ZZ[u]
\)，
并由命题 \ref{prop:witt-frobenius-defect-divisibility} 得到整除深度
\(
\Delta_{p,k}(u)\in p^k\ZZ[u]
\)。
另一方面，系数环版本的 Big-Witt/Möbius 反演在带 Adams 操作族 \(\psi^d\) 的交换环 \(R\) 上满足
\(
n\,p_n=\sum_{d\mid n}\mu(d)\,\psi^d(a_{n/d})\in R
\)
（命题 \ref{prop:witt-necklace-dwork-R}），其中在多参数多项式环情形 \(\psi^d\) 由逐变量幂替换实现。
将上述两条与端点三态压缩并置，可定义一个可直接审计实现链路的整数证书族。

\begin{definition}[端点三态上的 Frobenius 整数证书]\label{def:endpoint-frobenius-integer-certificate}
对任意素数 \(p\)、整数 \(k\ge 1\) 与端点态 \(\sigma\in\mathcal S_{\star}\)，取 \(r_\sigma\in\{\,\mathrm{i},1,-1\,\}\) 使 \(r_\sigma+r_\sigma^{-1}=\sigma\)，并令 \(u_\sigma:=r_\sigma^{2}\in\{-1,1\}\)。
定义
\begin{equation}\label{eq:endpoint-frobenius-integer-certificate}
\boxed{\ \mathfrak I_{p,k}(\sigma):=\frac{\Delta_{p,k}(u_{\sigma})}{p^{k}}\in\ZZ\ }.
\end{equation}
\end{definition}

\begin{proposition}[整数层先验护栏：端点证书的可审计性]\label{prop:endpoint-frobenius-integer-certificate-auditable}
对任意 \((p,k)\) 与 \(\sigma\in\mathcal S_{\star}\)，证书 \(\mathfrak I_{p,k}(\sigma)\) 为良定义的整数。
若在任一实现中出现 \(\mathfrak I_{p,k}(\sigma)\notin\ZZ\)（等价于 \(\Delta_{p,k}(u_\sigma)\notin p^k\ZZ\)），则必违反 Frobenius 缺陷整除深度 \(\Delta_{p,k}(u)\in p^k\ZZ[u]\)，从而可在整数层先行定位错误，再进入单位圆谱半径等连续谱层讨论。
\end{proposition}
\begin{proof}
由命题 \ref{prop:witt-frobenius-defect-divisibility}，\(\Delta_{p,k}(u)\) 的各系数都被 \(p^k\) 整除；因此对任意 \(u_\sigma\in\ZZ\) 有 \(\Delta_{p,k}(u_\sigma)\in p^k\ZZ\)，从而 \(\mathfrak I_{p,k}(\sigma)\in\ZZ\)。
若出现反例，则与“系数级整除”矛盾。证毕。
\end{proof}

\leanverified{paper\_endpoint\_frobenius\_integer\_certificate\_auditable}

\paragraph{整数层残差与谱层偏移的正交性：Frobenius 缺陷作为算术精度轴}
由 \(\Delta_{p,k}(u)\in p^k\ZZ[u]\) 可在整数层定义规范化残差多项式
\begin{equation}\label{eq:frobenius-normalized-residual-poly}
\boxed{\ E_{p,k}(u):=\frac{\Delta_{p,k}(u)}{p^{k}}\in\ZZ[u]\ }.
\end{equation}
于是 \(\{E_{p,k}\}_{p,k}\) 构成一个完全离散的“算术精度轴”：其审计只依赖整数系数与整除深度，而与单位圆谱半径、端点精度势等连续谱层指标无关。
在端点三态 \(\sigma\in\mathcal S_{\star}\) 上，上述残差与定义 \ref{def:endpoint-frobenius-integer-certificate} 的证书严格一致：
\[
\boxed{\ \mathfrak I_{p,k}(\sigma)=E_{p,k}(u_\sigma)\in\ZZ\ }.
\]
从而“先整数层、后谱层”的分层审计原则可被表述为：先以 \eqref{eq:frobenius-normalized-residual-poly}（或其端点取值）锁定实现的一致性护栏，
再进入端点精度势 \(y=-\log|\mathfrak{w}|\)（式 \eqref{eq:endpoint-busemann-potential}）与 unitary slice 上的连续谱偏移等分析层讨论。
